\section*{Chapitre \ref{ch:g12:curly-arrows} --- Mécanismes réactionnels~: les flèches courbes}

\begin{solution}{exo:g12:curly-arrows:1}
\ce{OH-}~: donneur (\omterm{def:g10:lewis-and-shape:lone-pair}{doublets non liants}, charge négative)~; aucun site
accepteur. \ce{H+}~: accepteur (aucun \omterm{def:g9:inside-the-atom:nucleus}{électron}). \ce{CH2=CH2}~: donneur
(la \omterm{def:g10:lewis-and-shape:multiple-bond}{double liaison}). \ce{CH3-Cl}~: donneur (\omterm{def:g10:lewis-and-shape:lone-pair}{doublets non liants} de
\ce{Cl}), accepteur (le carbone,
$\delta^+$). \ce{H2O}~: donneur (\omterm{def:g10:lewis-and-shape:lone-pair}{doublets non liants} de \ce{O})~; ses
\omterm{def:g7:atoms-and-molecules:atom}{atomes} d'hydrogène,
$\delta^+$, sont des sites accepteurs faibles.
\end{solution}

\begin{solution}{exo:g12:curly-arrows:2}
(a) \omterm{def:g12:curly-arrows:categories}{substitution}~; (b) \omterm{def:g12:curly-arrows:categories}{addition}~; (c) \omterm{def:g12:curly-arrows:categories}{élimination}~; (d) \omterm{def:g12:curly-arrows:categories}{addition}.
\end{solution}

\begin{solution}{exo:g12:curly-arrows:3}
Elle part d'un doublet d'\omterm{def:g9:inside-the-atom:nucleus}{électrons} (\omterm{def:g10:lewis-and-shape:lone-pair}{doublet non liant} ou liaison) et
aboutit sur l'\omterm{def:g7:atoms-and-molecules:atom}{atome} vers lequel va le doublet. D'un \omterm{def:g10:lewis-and-shape:lone-pair}{doublet non liant} de
l'oxygène vers un carbone~: le doublet devient une nouvelle liaison
\ce{O-C}.
\end{solution}

\begin{solution}{exo:g12:curly-arrows:4}
L'ammoniac cède le \omterm{def:g10:lewis-and-shape:lone-pair}{doublet non liant} de son azote~; la flèche va de ce
doublet vers \ce{H+}. Le produit \ce{NH4+} porte la charge $+1$.
\end{solution}

\begin{solution}{exo:g12:curly-arrows:5}
Une espèce formée lors d'une étape et consommée lors d'une étape
ultérieure~; elle ne figure pas dans l'équation globale. Ici, le
carbocation \ce{CH3-CH2+}.
\end{solution}

\begin{solution}{exo:g12:curly-arrows:6}
Une flèche d'un \omterm{def:g10:lewis-and-shape:lone-pair}{doublet non liant} de l'oxygène de \ce{OH-} vers le
carbone~; une flèche de la liaison \ce{C-Br} vers le brome. Charges~:
$-1$ à gauche
(\ce{OH-})~; à droite, le méthanol est neutre et \ce{Br-} porte $-1$~:
la charge totale vaut $-1$ des deux côtés.
\end{solution}

\begin{solution}{exo:g12:curly-arrows:7}
Le brome est plus électronégatif que l'hydrogène~: il prend le doublet
et part sous forme de \ce{Br-} ($-1$)~; l'hydrogène se lie à un carbone,
et l'autre carbone se retrouve avec une charge positive~:
\ce{CH3-CH2+} ($+1$).
\end{solution}

\begin{solution}{exo:g12:curly-arrows:8}
1.~\ce{CH3-CH2-OH + H+ -> CH3-CH2-OH2+}~: donneur, un \omterm{def:g10:lewis-and-shape:lone-pair}{doublet non liant}
de l'oxygène~; accepteur, \ce{H+}.
2.~\ce{CH3-CH2-OH2+ -> CH3-CH2+ + H2O}~: le doublet
\ce{C-O} va vers l'oxygène positif (accepteur), qui part avec lui sous
forme d'eau. 3.~\ce{CH3-CH2+ -> CH2=CH2 + H+}~: le doublet d'une
liaison \ce{C-H} (donneur) se déplace vers le carbone positif
(accepteur) pour former la \omterm{def:g10:lewis-and-shape:multiple-bond}{double liaison}, et \ce{H+} part. L'\omterm{def:g9:ions:ion}{ion}
hydrogène utilisé à l'étape 1 est restitué à l'étape 3~: il apparaît des
deux côtés et s'élimine.
\end{solution}

\begin{solution}{exo:g12:curly-arrows:9}
Une flèche d'un \omterm{def:g10:lewis-and-shape:lone-pair}{doublet non liant} de l'oxygène de l'eau vers le carbone
positif. Après perte de \ce{H+}, on obtient l'éthanol
\ce{CH3-CH2-OH}. Au total~:
\ce{CH2=CH2 + H2O -> CH3-CH2-OH}, une \omterm{def:g12:curly-arrows:categories}{addition}.
\end{solution}

\begin{solution}{exo:g12:curly-arrows:10}
Deux doublets se déplacent, en une seule étape~: la liaison \ce{C-O} se
forme, la liaison \ce{C-Br} se rompt. Pas d'intermédiaire.
\end{solution}

\begin{solution}{exo:g12:curly-arrows:11}
Le chlore est plus électronégatif ($3.16 - 2.55 = 0.61$ contre
$2.96 - 2.55 = 0.41$)~: le carbone du chlorométhane est le plus
$\delta^+$, ce qui, à lui seul, le ferait réagir plus vite. Ce n'est pas
le cas~: la vitesse dépend aussi de la facilité avec laquelle l'\omterm{def:g10:electron-shells:family}{halogène}
part avec le doublet, et le brome, plus gros, part plus facilement.
\end{solution}

\begin{solution}{exo:g12:curly-arrows:12}
Si l'hydrogène se lie au carbone 1~: \ce{CH3-CH+-CH3}, qui donne le
2-chloropropane \ce{CH3-CHCl-CH3}. S'il se lie au carbone 2~:
\ce{+CH2-CH2-CH3}, qui donne le 1-chloropropane \ce{CH2Cl-CH2-CH3}.
\end{solution}

\begin{solution}{exo:g12:curly-arrows:13}
La flèche est dans le mauvais sens~: elle doit partir du donneur, un
\omterm{def:g10:lewis-and-shape:lone-pair}{doublet non liant} de \ce{OH-}, et aboutir sur le carbone accepteur. Et
la nouvelle liaison se forme entre l'oxygène et le carbone, pas entre
l'oxygène et le brome~: le brome part en emportant le doublet de la
liaison \ce{C-Br}.
\end{solution}

\begin{solution}{exo:g12:curly-arrows:14}
Le carbone du groupe \ce{C=O} d'un acide est lié à deux \omterm{def:g7:atoms-and-molecules:atom}{atomes}
d'oxygène électronégatifs~: il est fortement $\delta^+$, c'est un
accepteur. L'oxygène de l'\omterm{def:g11:functional-groups:alcohol}{alcool} porte $\delta^-$ et deux \omterm{def:g10:lewis-and-shape:lone-pair}{doublets non
liants}~: c'est un donneur. De l'eau est perdue au total~; le \ce{-OH} de
l'acide est remplacé par le \ce{-O-R} de l'\omterm{def:g11:functional-groups:alcohol}{alcool}~: c'est une
\omterm{def:g12:curly-arrows:categories}{substitution}.
\end{solution}

\begin{solution}{exo:g12:curly-arrows:15}
Repoussé par les \omterm{def:g9:inside-the-atom:nucleus}{électrons} de la \omterm{def:g10:lewis-and-shape:multiple-bond}{double liaison}, le doublet de la
liaison \ce{Br-Br} se déplace vers le brome éloigné~: le brome proche
devient $\delta^+$, un site accepteur. Première étape~: une flèche de la
\omterm{def:g10:lewis-and-shape:multiple-bond}{double liaison} vers le brome proche, et une de la liaison \ce{Br-Br}
vers le brome éloigné, qui part sous forme de \ce{Br-}~; il reste une
charge positive sur l'autre carbone de l'ancienne \omterm{def:g10:lewis-and-shape:multiple-bond}{double liaison}.
\end{solution}

\begin{solution}{pb:g12:curly-arrows:1}
\textbf{1.} Éthène~: \ce{H2C=CH2}, une \omterm{def:g10:lewis-and-shape:multiple-bond}{double liaison} et quatre
\ce{C-H}. Bromure d'hydrogène~: \ce{H-Br} avec trois \omterm{def:g10:lewis-and-shape:lone-pair}{doublets non liants}
sur le brome.

\textbf{2.} La \omterm{def:g10:lewis-and-shape:multiple-bond}{double liaison}~: son second doublet constitue une région
riche en \omterm{def:g9:inside-the-atom:nucleus}{électrons}.

\textbf{3.} $2.96 - 2.20 = 0.76$~: \ce{H} porte $\delta^+$.

\textbf{4.} L'\omterm{def:g7:atoms-and-molecules:atom}{atome} d'hydrogène.

\textbf{5.} Une flèche de la \omterm{def:g10:lewis-and-shape:multiple-bond}{double liaison} vers l'hydrogène de
\ce{HBr}~; une flèche de la liaison \ce{H-Br} vers le brome.

\textbf{6.} Le carbocation \ce{CH3-CH2+} ($+1$) et \ce{Br-} ($-1$).

\textbf{7.} Avant~: $0 + 0 = 0$~; après~: $+1 - 1 = 0$.

\textbf{8.} Une flèche d'un \omterm{def:g10:lewis-and-shape:lone-pair}{doublet non liant} de \ce{Br-} vers le
carbone positif.

\textbf{9.} Six (trois liaisons)~: il lui manque un doublet pour avoir
un octet, il accepte donc un doublet de n'importe quel donneur voisin.

\textbf{10.} Le carbocation \ce{CH3-CH2+}.

\textbf{11.} Il est formé lors de la première étape et consommé lors de
la seconde.

\textbf{12.} Une \omterm{def:g12:curly-arrows:categories}{addition}.

\textbf{13.} Seulement le groupe~: la chaîne de deux carbones est
conservée, et la \omterm{def:g10:lewis-and-shape:multiple-bond}{double liaison} devient une liaison simple portant
\ce{H} et \ce{Br}.

\textbf{14.} \qty{28.0}{g/mol}; \qty{108.9}{g/mol}.

\textbf{15.} $10.0 / 28.0 = \qty{0.357}{mol}$.

\textbf{16.} Éthène~: $0.357 - x$~; \ce{HBr}~: en excès~; bromoéthane~:
$x$. L'éthène est limitant~: $x_{\max} = \qty{0.357}{mol}$.

\textbf{17.} $0.357 \times 108.9 = \qty{38.9}{g}$.

\textbf{18.} $0.75 \times 38.9 = \qty{29.2}{g}$.

\textbf{19.} Environ \qty{29}{g} de bromoéthane.
\end{solution}
